% =============================================================================
% LangChain Course Resumen / Code Explainer
% Based on the local project modules 01–05
% Compile with: pdflatex langchain-resumen.tex  (run twice for TOC)
% =============================================================================
\documentclass[11pt,a4paper]{article}

\usepackage[margin=2.2cm]{geometry}
\usepackage[T1]{fontenc}
\usepackage[utf8]{inputenc}
\usepackage{lmodern}
\usepackage{microtype}
\usepackage{enumitem}
\usepackage{booktabs}
\usepackage{longtable}
\usepackage{tabularx}
\usepackage{xcolor}
\usepackage{listings}
\usepackage{hyperref}
\usepackage{titlesec}
\usepackage{fancyhdr}
\usepackage{tcolorbox}
\usepackage{graphicx}
\usepackage{amsmath}
\usepackage{multirow}
\usepackage{array}

\hypersetup{
  colorlinks=true,
  linkcolor=blue!60!black,
  urlcolor=blue!50!black,
  citecolor=blue!60!black
}

% --- Code listing style ---
\definecolor{codebg}{RGB}{248,248,248}
\definecolor{codeframe}{RGB}{200,200,200}
\definecolor{kwcolor}{RGB}{0,80,140}
\definecolor{strcolor}{RGB}{140,40,40}
\definecolor{cmtcolor}{RGB}{100,120,100}

\lstdefinestyle{py}{
  language=Python,
  backgroundcolor=\color{codebg},
  basicstyle=\ttfamily\footnotesize,
  keywordstyle=\color{kwcolor}\bfseries,
  stringstyle=\color{strcolor},
  commentstyle=\color{cmtcolor}\itshape,
  breaklines=true,
  frame=single,
  rulecolor=\color{codeframe},
  showstringspaces=false,
  tabsize=4,
  columns=fullflexible,
  keepspaces=true,
  % listings does not handle UTF-8 natively; map Spanish chars
  literate=
    {á}{{\'a}}1 {é}{{\'e}}1 {í}{{\'i}}1 {ó}{{\'o}}1 {ú}{{\'u}}1
    {Á}{{\'A}}1 {É}{{\'E}}1 {Í}{{\'I}}1 {Ó}{{\'O}}1 {Ú}{{\'U}}1
    {ñ}{{\~n}}1 {Ñ}{{\~N}}1 {ü}{{\"{u}}}1 {Ü}{{\"{U}}}1
}
\lstset{style=py}

% --- Boxes ---
\tcbuselibrary{skins,breakable}
\newtcolorbox{concept}[1][]{
  colback=blue!4, colframe=blue!50!black, fonttitle=\bfseries,
  title={Concepto}, breakable, #1
}
\newtcolorbox{tip}[1][]{
  colback=green!4, colframe=green!50!black, fonttitle=\bfseries,
  title={Tip / recuerdo}, breakable, #1
}
\newtcolorbox{warn}[1][]{
  colback=orange!6, colframe=orange!70!black, fonttitle=\bfseries,
  title={Cuidado}, breakable, #1
}
\newtcolorbox{fromcode}[1][]{
  colback=gray!6, colframe=gray!50, fonttitle=\bfseries\ttfamily,
  title={Del proyecto}, breakable, #1
}

% --- Header ---
\pagestyle{fancy}
\fancyhf{}
\fancyhead[L]{\small LangChain Course --- Resumen}
\fancyhead[R]{\small Modules 01--05}
\fancyfoot[C]{\thepage}
\renewcommand{\headrulewidth}{0.4pt}

\title{%
  \textbf{LangChain --- Resumen \& Code Explainer}\\[0.4em]
  \large Hoja de estudio a partir del curso local
}
\author{Basado en \texttt{langchain-course} (módulos 01--05)}
\date{}

\begin{document}
\maketitle
\tableofcontents
\vspace{1em}
\noindent\textit{Propósito:} hoja de referencia para recordar cómo funciona LangChain en este proyecto: mensajes, tools, agents, LCEL (\texttt{\textbar}), RAG, Pinecone/Chroma, Tavily, Ollama, API keys y debugging con LangSmith.

% =============================================================================
\section{Mapa del curso (qué hace cada carpeta)}
% =============================================================================

\begin{center}
\begin{tabular}{@{}clp{8.2cm}@{}}
\toprule
\textbf{Módulo} & \textbf{Carpeta} & \textbf{Idea central} \\
\midrule
01 & \texttt{01-hello-world} & Setup: \texttt{dotenv}, punto de entrada \\
02 & \texttt{02-web-search-tool} & Agent + tools + Tavily + output estructurado \\
03 & \texttt{03-agents-under-the-hood} & ReAct manual vs tool-calling nativo \\
04 & \texttt{04-rag} & Ingestión + retrieval + LCEL vs imperativo \\
05 & \texttt{05-building-\ldots} & Crawl docs (Tavily) + RAG agent + Streamlit \\
\bottomrule
\end{tabular}
\end{center}

\begin{concept}[title=¿Qué es LangChain?]
LangChain es un \textbf{framework de orquestación} para apps con LLMs.
No es el modelo: conecta modelos, prompts, tools, retrievers y memoria en \textbf{pipelines} (chains) o \textbf{agents} (bucles que deciden qué tool usar).
\end{concept}

% =============================================================================
\section{Setup: entorno, paquetes y API keys}
% =============================================================================

\subsection{Gestión del proyecto}
Este repo usa \textbf{uv} + \texttt{pyproject.toml}. Dependencias clave (resumidas):

\begin{itemize}[leftmargin=*,itemsep=2pt]
  \item \texttt{langchain}, \texttt{langchain-community}, \texttt{langchain-core} --- núcleo
  \item \texttt{langchain-ollama}, \texttt{langchain-google-genai}, \texttt{langchain-openai} --- modelos
  \item \texttt{langchain-pinecone}, \texttt{langchain-chroma} --- vector stores
  \item \texttt{langchain-tavily}, \texttt{tavily-python} --- search / crawl
  \item \texttt{langchain-text-splitters}, \texttt{langchain-unstructured} --- chunking / loaders
  \item \texttt{langsmith} --- tracing / debugging
  \item \texttt{python-dotenv}, \texttt{streamlit}, \texttt{pydantic}
\end{itemize}

\subsection{Variables de entorno (\texttt{.env})}
Siempre al inicio de los scripts:

\begin{lstlisting}
from dotenv import load_dotenv
load_dotenv()  # lee .env y rellena os.environ
\end{lstlisting}

Variables típicas en este curso (nombres, \textbf{nunca} valores reales en notas):

\begin{center}
\begin{tabular}{@{}llp{6.5cm}@{}}
\toprule
\textbf{Variable} & \textbf{Servicio} & \textbf{Para qué} \\
\midrule
\texttt{OPENAI\_API\_KEY} & OpenAI & Chat + embeddings cloud \\
\texttt{GOOGLE\_API\_KEY} & Gemini & \texttt{ChatGoogleGenerativeAI} \\
\texttt{PINECONE\_API\_KEY} & Pinecone & Vector DB en la nube \\
\texttt{INDEX\_NAME} & Pinecone & Nombre del índice \\
\texttt{TAVILY\_API\_KEY} & Tavily & Search / Map / Extract / Crawl \\
\texttt{LANGCHAIN\_API\_KEY} & LangSmith & Tracing \\
\texttt{LANGCHAIN\_TRACING\_V2} & LangSmith & \texttt{true} activa traces \\
\texttt{LANGCHAIN\_PROJECT} & LangSmith & Nombre del proyecto en UI \\
\texttt{LANGCHAIN\_ENDPOINT} & LangSmith & API (p.ej. región EU) \\
\bottomrule
\end{tabular}
\end{center}

\begin{tip}[title=Modelos locales vs cloud]
Con \textbf{Ollama} no necesitas API key del LLM: el modelo corre en tu máquina
(\texttt{ChatOllama}, \texttt{OllamaEmbeddings}).
Cloud (OpenAI, Gemini) sí requiere key en \texttt{.env}.
Pinecone y Tavily siempre necesitan sus keys si los usas.
\end{tip}

\begin{warn}
Nunca subas \texttt{.env} a git. Las keys dan acceso de facturación.
Si una key se filtra, rotarla en el dashboard del proveedor.
\end{warn}

% =============================================================================
\section{Arquitectura mental: Chain vs Agent}
% =============================================================================

\begin{center}
\begin{tabular}{@{}p{7.2cm}p{7.2cm}@{}}
\toprule
\textbf{Chain (pipeline fijo)} & \textbf{Agent (bucle de decisión)} \\
\midrule
Pasos definidos por ti:
\texttt{retrieve $\to$ prompt $\to$ llm $\to$ parse}
& El LLM elige tools y cuántas veces llamarlas \\
Ejemplo: RAG con LCEL en \texttt{04-rag/main.py}
& Ejemplo: \texttt{create\_agent} en 02 y 05;
bucle manual en 03 \\
Más predecible y barato
& Más flexible, más tokens / latencia \\
\bottomrule
\end{tabular}
\end{center}

% =============================================================================
\section{Mensajes: el ``idioma'' de LangChain}
% =============================================================================

Los chat models no reciben un string suelto idealmente: reciben una \textbf{lista de mensajes} tipados.

\begin{center}
\begin{tabular}{@{}llp{7.5cm}@{}}
\toprule
\textbf{Clase} & \textbf{Rol} & \textbf{Uso} \\
\midrule
\texttt{SystemMessage} & system & Instrucciones / reglas del asistente \\
\texttt{HumanMessage} & user & Pregunta del usuario \\
\texttt{AIMessage} & assistant & Respuesta del modelo (puede incluir \texttt{tool\_calls}) \\
\texttt{ToolMessage} & tool & Resultado de ejecutar una tool (lleva \texttt{tool\_call\_id}) \\
\bottomrule
\end{tabular}
\end{center}

\begin{fromcode}[title=03-agents-under-the-hood/agent-loop-tool-calling.py]
\begin{lstlisting}
messages = [
    SystemMessage(content="You are a helpful assistant..."),
    HumanMessage(content=question),
]
ai_message = llm_with_tools.invoke(messages)
# si hay tool_calls -> ejecutas tool -> append AIMessage + ToolMessage
messages.append(ai_message)
messages.append(ToolMessage(content=str(observation),
                            tool_call_id=tool_call_id))
\end{lstlisting}
\end{fromcode}

\begin{concept}[title=¿Por qué ToolMessage necesita tool\_call\_id?]
El modelo puede pedir varias tools. El \texttt{tool\_call\_id} enlaza cada resultado
con la petición concreta. Sin eso, el historial queda ambiguo.
\end{concept}

Forma alternativa (dicts), usada en el doc assistant:

\begin{lstlisting}
messages = [{"role": "user", "content": query}]
\end{lstlisting}

Ambas son válidas; las clases dan tipado y helpers de LangChain.

% =============================================================================
\section{\texttt{invoke}, \texttt{stream}, \texttt{batch}, async}
% =============================================================================

Casi todo objeto ``Runnable'' en LangChain (LLM, chain, retriever, tool, agent) expone:

\begin{center}
\begin{tabular}{@{}llp{7.8cm}@{}}
\toprule
\textbf{Método} & \textbf{Sync/Async} & \textbf{Significado} \\
\midrule
\texttt{invoke(input)} & sync & Una entrada $\to$ una salida completa \\
\texttt{stream(input)} & sync & Tokens / chunks conforme se generan \\
\texttt{batch([..])} & sync & Varias entradas en paralelo (batch) \\
\texttt{ainvoke} & async & Igual que invoke, no bloquea el event loop \\
\texttt{astream} & async & Streaming asíncrono \\
\texttt{aadd\_documents} & async & Indexar docs en vector store (05) \\
\bottomrule
\end{tabular}
\end{center}

\begin{tip}
Si ves \texttt{.invoke(...)}, estás ejecutando un Runnable.
El \textbf{tipo del input} depende del objeto:
LLM $\to$ mensajes o string; chain LCEL $\to$ dict (\texttt{\{"question": \ldots\}});
agent $\to$ \texttt{\{"messages": [\ldots]\}}.
\end{tip}

% =============================================================================
\section{Modelos de chat (LLM wrappers)}
% =============================================================================

\subsection{Formas usadas en el curso}

\begin{lstlisting}
# Local (Ollama)
from langchain_ollama import ChatOllama
llm = ChatOllama(model="qwen3", temperature=0)

# Cloud Gemini
from langchain_google_genai import ChatGoogleGenerativeAI
llm = ChatGoogleGenerativeAI(model="gemini-...", temperature=0.2)

# Factory generica (elige provider por prefijo / config)
from langchain.chat_models import init_chat_model
llm = init_chat_model("ollama:qwen3:1.7b", temperature=0)
llm = init_chat_model("gpt-5.2", model_provider="openai")
\end{lstlisting}

Parámetros importantes:
\begin{itemize}[leftmargin=*]
  \item \texttt{temperature}: 0 = más determinista (útil en agents/tools); $>$0 = más creativo.
  \item El \textbf{nombre del modelo} debe existir en Ollama (\texttt{ollama pull ...}) o en el proveedor cloud.
\end{itemize}

\subsection{\texttt{bind\_tools}: conectar tools al modelo}
\begin{lstlisting}
llm_with_tools = llm.bind_tools(tools)
ai_message = llm_with_tools.invoke(messages)
tool_calls = ai_message.tool_calls  # lista de dicts: name, args, id
\end{lstlisting}

Esto es \textbf{native function calling}: el modelo responde con una estructura de tool call,
no con texto ``Action: \ldots'' que tú parseas a mano.

% =============================================================================
\section{Tools: el decorador \texttt{@tool}}
% =============================================================================

\begin{lstlisting}
from langchain.tools import tool

@tool
def get_product_price(product_name: str) -> float:
    """Look up the price of a product in the catalogue.

    Args:
        product_name: The base name of the product ONLY (e.g., 'laptop').
    """
    return prices.get(product_name, 0)
\end{lstlisting}

\begin{concept}[title=Qué hace @tool]
Convierte una función Python en una \textbf{Tool} de LangChain:
\begin{itemize}[leftmargin=*]
  \item Nombre = nombre de la función
  \item Descripción = docstring (el LLM la lee para decidir cuándo usarla)
  \item Schema de args = type hints (+ Pydantic si hace falta)
  \item Ejecutable con \texttt{tool.invoke(args\_dict)}
\end{itemize}
\textbf{Docstring bueno = tool que el modelo usa bien.} Si el docstring es vago, el agent falla.
\end{concept}

Variante avanzada (doc assistant): devolver contenido \emph{y} artefacto:

\begin{lstlisting}
@tool(response_format="content_and_artifact")
def retrieve_context(query: str):
    """Retrieve relevant documentation..."""
    docs = vectorstore.as_retriever().invoke(query, k=4)
    serialized = "..."  # texto para el LLM
    return serialized, retrieved_docs  # artifact para la UI (fuentes)
\end{lstlisting}

El LLM ve el texto; tu app puede leer \texttt{ToolMessage.artifact} (lista de \texttt{Document}) para citar fuentes.

% =============================================================================
\section{Agents: tres niveles de abstracción}
% =============================================================================

\subsection{Nivel A --- ReAct ``a mano'' (prompt + regex)}
Archivo: \texttt{03-agents-under-the-hood/react-prompt-function-calling.py}

\begin{itemize}[leftmargin=*]
  \item Prompt estilo ReAct: Thought / Action / Action Input / Observation
  \item El LLM escribe texto; tú parseas con regex
  \item Tools son funciones normales (aquí con \texttt{@traceable})
  \item \texttt{stop=["\textbackslash nObservation"]} fuerza al modelo a parar antes de inventar la observación
\end{itemize}

\begin{tip}[title=Para qué sirve este módulo]
Entender \textbf{qué hace un agent por debajo}. En producción preferirás native tool calling
(Nivel B/C), más robusto que parsear texto.
\end{tip}

\subsection{Nivel B --- Agent loop con tool calling nativo}
Archivo: \texttt{agent-loop-tool-calling.py}

Flujo por iteración:
\begin{enumerate}[leftmargin=*]
  \item \texttt{llm\_with\_tools.invoke(messages)}
  \item Si \textbf{no} hay \texttt{tool\_calls} $\to$ respuesta final
  \item Si hay $\to$ ejecutar tool $\to$ añadir \texttt{AIMessage} + \texttt{ToolMessage}
  \item Repetir hasta \texttt{MAX\_ITERATIONS}
\end{enumerate}

En el código del curso se fuerza \textbf{una tool por iteración} (solo \texttt{tool\_calls[0]}),
para control didáctico del orden: precio $\to$ descuento $\to$ shipping.

\subsection{Nivel C --- \texttt{create\_agent} (API moderna)}
Archivos: \texttt{02-web-search-tool/web\_search.py}, \texttt{05/\ldots/backend/core.py}

\begin{lstlisting}
from langchain.agents import create_agent

agent = create_agent(
    model=llm,
    tools=tools,
    response_format=AgentResponse,  # opcional: salida Pydantic
    # system_prompt="..."           # en el doc helper
)

result = agent.invoke({
    "messages": [HumanMessage(content="...")]
})
\end{lstlisting}

\begin{concept}[title=create\_agent]
Abstrae el bucle: tool binding, mensajes, iteraciones y (opcional) formato de respuesta.
Tú defines tools + prompt; LangChain gestiona el loop.
\end{concept}

\subsection{Salida estructurada con Pydantic}
\begin{lstlisting}
class Source(BaseModel):
    url: str = Field(description="The URL of the source")

class AgentResponse(BaseModel):
    answer: str = Field(description="The answer...")
    sources: List[Source] = Field(description="Sources used...")

agent = create_agent(model=llm, tools=tools, response_format=AgentResponse)
\end{lstlisting}

El modelo debe devolver JSON conforme al schema --- útil para UIs y APIs tipadas.

% =============================================================================
\section{LCEL y el operador pipe \texttt{|}}
% =============================================================================

\textbf{LCEL} = LangChain Expression Language.
Los Runnables se componen con \texttt{|} como un pipeline Unix.

\subsection{Ejemplo del curso (\texttt{04-rag/main.py})}

\begin{lstlisting}
from operator import itemgetter
from langchain_core.runnables import RunnablePassthrough
from langchain_core.output_parsers import StrOutputParser

retrieval_chain = (
    RunnablePassthrough.assign(
        context=itemgetter("question") | retriever | format_docs
    )
    | prompt_template
    | llm
    | StrOutputParser()
)

answer = retrieval_chain.invoke({"question": query})
\end{lstlisting}

Lectura de izquierda a derecha:
\begin{enumerate}[leftmargin=*]
  \item Entrada: \texttt{\{"question": "..."\}}
  \item \texttt{RunnablePassthrough.assign(...)} copia el dict y \textbf{añade} la clave \texttt{context}:
        \begin{itemize}
          \item saca \texttt{question} con \texttt{itemgetter}
          \item la pasa al \texttt{retriever} $\to$ lista de \texttt{Document}
          \item \texttt{format\_docs} las une en un string
        \end{itemize}
  \item \texttt{prompt\_template} rellena \texttt{\{context\}} y \texttt{\{question\}}
  \item \texttt{llm} genera respuesta
  \item \texttt{StrOutputParser()} deja solo el string (\texttt{.content})
\end{enumerate}

\begin{center}
\begin{tabular}{@{}p{7.2cm}p{7.2cm}@{}}
\toprule
\textbf{Sin LCEL (imperativo)} & \textbf{Con LCEL} \\
\midrule
\texttt{docs = retriever.invoke(...)}\\
\texttt{context = format\_docs(docs)}\\
\texttt{messages = prompt.format\_...}\\
\texttt{response = llm.invoke(...)}
& Una expresión componible;
\texttt{stream}/\texttt{ainvoke}/\texttt{batch} gratis;
mejor observabilidad
\\
\bottomrule
\end{tabular}
\end{center}

\begin{tip}[title=Analogía]
\texttt{a | b | c} $\approx$ ``pasa la salida de a a b, luego a c''.
Cada pieza debe ser un Runnable (o función envuelta como tal).
\end{tip}

Piezas frecuentes de LCEL:
\begin{itemize}[leftmargin=*]
  \item \texttt{ChatPromptTemplate.from\_template("... \{var\} ...")}
  \item \texttt{RunnablePassthrough} / \texttt{.assign(...)}
  \item \texttt{itemgetter("clave")} --- del módulo \texttt{operator}
  \item \texttt{StrOutputParser()}
\end{itemize}

% =============================================================================
\section{RAG: Retrieval-Augmented Generation}
% =============================================================================

\begin{concept}[title=Idea de RAG]
El LLM no ``sabe'' tus PDFs/docs. RAG:
\begin{enumerate}[leftmargin=*]
  \item \textbf{Ingiere} documentos $\to$ chunks $\to$ embeddings $\to$ vector store
  \item En la pregunta: \textbf{recupera} chunks similares
  \item Pasa esos chunks como \textbf{contexto} al prompt + pregunta
  \item El LLM responde \emph{basado en el contexto} (menos alucinación)
\end{enumerate}
\end{concept}

\subsection{Pipeline de ingestión (patrón del curso)}

\begin{center}
\texttt{Loader} $\to$ \texttt{TextSplitter} $\to$ \texttt{Embeddings} $\to$ \texttt{VectorStore.from\_documents}
\end{center}

\subsubsection{Loaders}
\begin{lstlisting}
# Texto / unstructured
from langchain_unstructured import UnstructuredLoader
loader = UnstructuredLoader(file_path="...", chunking_strategy="basic")

# PDF
from langchain_community.document_loaders import PyPDFLoader
loader = PyPDFLoader("HuckFinn.pdf")
docs = loader.load()  # List[Document]
\end{lstlisting}

Un \texttt{Document} = \texttt{page\_content} (str) + \texttt{metadata} (dict, p.ej. \texttt{source}, página).

\subsubsection{Splitters (chunking)}
\begin{lstlisting}
from langchain_text_splitters import CharacterTextSplitter
from langchain_text_splitters import RecursiveCharacterTextSplitter

CharacterTextSplitter(chunk_size=1000, chunk_overlap=0)
RecursiveCharacterTextSplitter(chunk_size=500, chunk_overlap=50)
# En docs largas del crawl:
RecursiveCharacterTextSplitter(chunk_size=4000, chunk_overlap=200)
\end{lstlisting}

\begin{itemize}[leftmargin=*]
  \item \texttt{chunk\_size}: tamaño objetivo del trozo
  \item \texttt{chunk\_overlap}: solape para no cortar ideas a mitad
  \item \texttt{RecursiveCharacterTextSplitter}: intenta cortar por párrafos/frases (mejor calidad habitual)
\end{itemize}

\subsubsection{Embeddings}
Convierten texto $\to$ vector numérico (misma ``semántica'' $\approx$ vectores cercanos).

\begin{lstlisting}
from langchain_ollama import OllamaEmbeddings           # o community
from langchain_community.embeddings import OllamaEmbeddings
embeddings = OllamaEmbeddings(model="mxbai-embed-large")
# tambien: nomic-embed-text, text-embedding-3-small (OpenAI), etc.
\end{lstlisting}

\begin{warn}
El modelo de embedding al \textbf{consultar} debe ser el \textbf{mismo} (o compatible)
que al \textbf{ingestar}. Si no, la similitud no tiene sentido.
\end{warn}

\subsubsection{Vector stores: Pinecone vs Chroma}
\begin{center}
\begin{tabular}{@{}llp{6.8cm}@{}}
\toprule
 & \textbf{Pinecone} & \textbf{Chroma} \\
\midrule
Dónde & Nube & Local (\texttt{persist\_directory}) \\
Key & \texttt{PINECONE\_API\_KEY} & No \\
Uso en curso & 04-rag, backend core & 05 \texttt{ingestion.py} \\
API típica & \texttt{PineconeVectorStore} & \texttt{Chroma(...)} \\
\bottomrule
\end{tabular}
\end{center}

\begin{lstlisting}
# Escribir (ingest)
PineconeVectorStore.from_documents(texts, embeddings,
                                   index_name=os.environ["INDEX_NAME"])

# Leer (query)
vectorstore = PineconeVectorStore(index_name=..., embedding=embeddings)
retriever = vectorstore.as_retriever(search_kwargs={"k": 3})
docs = retriever.invoke(query)
\end{lstlisting}

\texttt{k} = cuántos vecinos más cercanos devolver (más contexto $\neq$ siempre mejor: ruido + tokens).

\subsubsection{Formatear contexto}
\begin{lstlisting}
def format_docs(docs):
    return "\n\n".join(doc.page_content for doc in docs)
\end{lstlisting}

\subsubsection{Prompt RAG típico}
\begin{lstlisting}
prompt = ChatPromptTemplate.from_template(
"""Answer based only on the following context:

{context}

Question: {question}

Provide a detailed answer:""")
\end{lstlisting}

La frase ``based only on the following context'' reduce inventar hechos fuera del retrieval.

\subsection{RAG ``simple'' vs RAG ``agentic''}
\begin{itemize}[leftmargin=*]
  \item \textbf{04-rag}: retrieval fijo $\to$ LLM (chain)
  \item \textbf{05 backend/core.py}: el agent tiene una tool \texttt{retrieve\_context};
        decide \emph{si} y \emph{con qué query} recuperar --- más flexible
\end{itemize}

% =============================================================================
\section{Tavily: search y crawl para agentes / RAG}
% =============================================================================

Tavily es una API orientada a LLM apps (resultados limpios, menos HTML basura).

\begin{center}
\begin{tabular}{@{}llp{7cm}@{}}
\toprule
\textbf{Clase} & \textbf{Paquete} & \textbf{Rol} \\
\midrule
\texttt{TavilySearch} & \texttt{langchain\_tavily} & Búsqueda web (tool de agent) \\
\texttt{TavilyMap} & idem & Descubre URLs de un sitio (sitemap crawl) \\
\texttt{TavilyExtract} & idem & Extrae contenido limpio de URLs \\
\texttt{TavilyCrawl} & idem & Crawl + extract en un paso \\
\bottomrule
\end{tabular}
\end{center}

Parámetros útiles de Map/Crawl:
\begin{itemize}[leftmargin=*]
  \item \texttt{max\_depth}: profundidad de enlaces
  \item \texttt{max\_breadth}: ramificación por nivel
  \item \texttt{max\_pages}: tope de páginas
  \item \texttt{extract\_depth}: p.ej. \texttt{"advanced"}
\end{itemize}

\begin{fromcode}[title=05 ingestion.py --- crawl $\to$ Documents]
\begin{lstlisting}
res = tavily_crawl.invoke({
    "url": "https://python.langchain.com/",
    "max_depth": 1,
    "extract_depth": "advanced",
})
for item in res["results"]:
    Document(page_content=item["raw_content"],
             metadata={"source": item["url"]})
\end{lstlisting}
\end{fromcode}

Flujo tutorial típico (notebooks):
\texttt{TavilyMap} (listar URLs) $\to$ \texttt{TavilyExtract} (contenido) $\to$ chunk $\to$ vector store.
\texttt{TavilyCrawl} acorta ese pipeline.

Requiere \texttt{TAVILY\_API\_KEY} en el entorno.

% =============================================================================
\section{Debugging y observabilidad (LangSmith)}
% =============================================================================

Con tracing activo, cada \texttt{invoke} aparece en la UI de LangSmith: prompts, tools, latencias, errores.

\begin{lstlisting}
# .env
LANGCHAIN_TRACING_V2=true
LANGCHAIN_API_KEY=...
LANGCHAIN_PROJECT=mi-proyecto
LANGCHAIN_ENDPOINT=https://eu.api.smith.langchain.com  # si usas region EU
\end{lstlisting}

Marcar funciones concretas:

\begin{lstlisting}
from langsmith import traceable

@traceable(name="LangChain Agent Loop")
def run_agent(question: str):
    ...

@traceable(run_type="tool")
def get_product_price(...):
    ...

@traceable(name="Ollama Chat Traced", run_type="llm")
def ollama_chat_traced(...):
    ...
\end{lstlisting}

\begin{tip}[title=Cómo depurar un agent]
\begin{enumerate}[leftmargin=*]
  \item Mira el trace: ¿llamó la tool correcta? ¿args raros?
  \item Revisa el \textbf{docstring} de la tool y el \textbf{SystemMessage}
  \item Baja \texttt{temperature} a 0
  \item Imprime \texttt{ai\_message.tool\_calls} y las observations (como en módulo 03)
  \item Limita \texttt{MAX\_ITERATIONS} para no gastar tokens en bucles
\end{enumerate}
\end{tip}

% =============================================================================
\section{UI: Streamlit en el Documentation Helper}
% =============================================================================

\texttt{05/.../main.py}: chat UI que llama a \texttt{run\_llm} del backend.

Ideas clave:
\begin{itemize}[leftmargin=*]
  \item \texttt{st.session\_state.messages} --- historial entre reruns
  \item \texttt{st.chat\_input} / \texttt{st.chat\_message} --- UI de chat
  \item Fuentes desde \texttt{Document.metadata["source"]}
  \item El backend puede usar Pinecone + OpenAI; la ingestión del curso también muestra Chroma + Ollama
\end{itemize}

% =============================================================================
\section{Glosario rápido (clases y paquetes)}
% =============================================================================

\begin{longtable}{@{}p{4.2cm}p{11cm}@{}}
\toprule
\textbf{Nombre} & \textbf{Qué es} \\
\midrule
\endhead
\texttt{langchain} & API de alto nivel (agents, chat\_models helpers) \\
\texttt{langchain\_core} & Mensajes, prompts, runnables, parsers, documents \\
\texttt{langchain\_community} & Integraciones community (loaders, algunos embeddings) \\
\texttt{ChatOllama} & Chat model local vía Ollama \\
\texttt{OllamaEmbeddings} & Embeddings locales \\
\texttt{init\_chat\_model} & Factory ``string $\to$ chat model'' \\
\texttt{create\_agent} & Construye agent con tools (+ system prompt / schema) \\
\texttt{@tool} & Decorador función $\to$ Tool \\
\texttt{bind\_tools} & Adjunta schemas de tools al LLM \\
\texttt{HumanMessage} & Mensaje del usuario \\
\texttt{SystemMessage} & Instrucciones de sistema \\
\texttt{AIMessage} & Respuesta del modelo (posible \texttt{tool\_calls}) \\
\texttt{ToolMessage} & Resultado de tool hacia el modelo \\
\texttt{ChatPromptTemplate} & Plantilla de prompt con variables \\
\texttt{RunnablePassthrough} & Pasa / enriquece el dict en LCEL \\
\texttt{StrOutputParser} & Extrae string de \texttt{AIMessage} \\
\texttt{Document} & Unidad de texto + metadata \\
\texttt{as\_retriever()} & Vector store $\to$ interfaz retrieve \\
\texttt{PineconeVectorStore} & Índice vectorial en Pinecone \\
\texttt{Chroma} & Vector store local persistente \\
\texttt{TavilySearch/Map/...} & Tools/APIs de Tavily \\
\texttt{traceable} & Decorador LangSmith para spans \\
\texttt{BaseModel} / \texttt{Field} & Schemas Pydantic (structured output) \\
\texttt{load\_dotenv} & Carga secrets desde \texttt{.env} \\
\bottomrule
\end{longtable}

% =============================================================================
\section{Cheatsheet de patrones del repo}
% =============================================================================

\subsection{1. Preguntar a un LLM (sin RAG)}
\begin{lstlisting}
result = llm.invoke([HumanMessage(content=query)])
print(result.content)
\end{lstlisting}

\subsection{2. Agent con tools}
\begin{lstlisting}
agent = create_agent(model=llm, tools=[my_tool])
result = agent.invoke({"messages": [HumanMessage(content="...")]})
\end{lstlisting}

\subsection{3. Ingest PDF \texorpdfstring{$\to$}{->} Pinecone}
\begin{lstlisting}
docs = PyPDFLoader(path).load()
chunks = RecursiveCharacterTextSplitter(
    chunk_size=500, chunk_overlap=50
).split_documents(docs)
PineconeVectorStore.from_documents(chunks, embeddings, index_name=INDEX)
\end{lstlisting}

\subsection{4. RAG chain LCEL}
\begin{lstlisting}
chain = (
    RunnablePassthrough.assign(
        context=itemgetter("question") | retriever | format_docs
    )
    | prompt | llm | StrOutputParser()
)
chain.invoke({"question": "..."})
\end{lstlisting}

\subsection{5. Retrieval como tool de agent}
\begin{lstlisting}
@tool(response_format="content_and_artifact")
def retrieve_context(query: str):
    docs = vectorstore.as_retriever().invoke(query, k=4)
    return serialize(docs), docs

agent = create_agent(model, tools=[retrieve_context], system_prompt=...)
\end{lstlisting}

\subsection{6. Orden típico de un agent ``calculadora de tools''}
En el ejemplo de precios (módulo 03):
\begin{enumerate}[leftmargin=*]
  \item \texttt{get\_product\_price}
  \item \texttt{apply\_discount}
  \item \texttt{calculate\_shipping}
\end{enumerate}
Las reglas van en el \texttt{SystemMessage} / ReAct prompt: el modelo \textbf{no} debe calcular a mano.

% =============================================================================
\section{Errores mentales frecuentes}
% =============================================================================

\begin{enumerate}[leftmargin=*]
  \item \textbf{Confundir LLM con embeddings}: uno genera texto; el otro vectores para búsqueda.
  \item \textbf{Olvidar \texttt{load\_dotenv()}}: ``missing API key'' aunque el \texttt{.env} exista.
  \item \textbf{Índice Pinecone vacío}: preguntar sin haber corrido ingestión.
  \item \textbf{Embedding distinto} en ingest vs query.
  \item \textbf{Docstring pobre} en \texttt{@tool} $\to$ el agent llama mal o no llama.
  \item \textbf{Esperar que RAG cite solo}: debes pedir citas en el system prompt y/o pasar \texttt{metadata.source}.
  \item \textbf{ReAct texto vs tool calling}: parsear ``Action:'' es frágil; \texttt{bind\_tools} es el camino moderno.
  \item \textbf{Input shape del agent}: casi siempre \lstinline|{"messages": [...]}|, no un string suelto.
\end{enumerate}

% =============================================================================
\section{Orden de lectura recomendado (repaso)}
% =============================================================================

\begin{enumerate}[leftmargin=*]
  \item \texttt{01-hello-world} --- dotenv y entrada
  \item \texttt{02-web-search-tool} --- \texttt{create\_agent}, \texttt{HumanMessage}, structured output
  \item \texttt{03/.../agent-loop-tool-calling.py} --- mensajes + \texttt{bind\_tools} + loop
  \item \texttt{03/.../react-prompt-...py} --- ReAct clásico (contexto histórico)
  \item \texttt{04-rag/ingestion.py} + \texttt{ingest-pdf.py} --- loader, split, embed, Pinecone
  \item \texttt{04-rag/main.py} --- RAG sin/con LCEL y el pipe \texttt{|}
  \item \texttt{05/ingestion.py} --- TavilyCrawl + Chroma/async batches
  \item \texttt{05/backend/core.py} + \texttt{main.py} --- RAG agentic + UI
\end{enumerate}

\vspace{1.5em}
\begin{center}
\fbox{\parbox{0.92\textwidth}{\centering
\textbf{Frase resumen}\\[0.4em]
\textit{LangChain conecta modelos, tools y datos.\\
Las chains componen pasos fijos con \texttt{|}.\\
Los agents dejan que el LLM elija tools en un bucle de mensajes.\\
RAG mete documentos relevantes en ese bucle o pipeline.\\
LangSmith te deja ver cada paso.}
}}
\end{center}

\end{document}
